\documentclass[11pt]{article}
\usepackage[a4paper,margin=0.65in]{geometry}
\usepackage[T1]{fontenc}
\usepackage{amsmath,amssymb,amsfonts}
\usepackage{graphicx}
\usepackage{pdfpages}
\usepackage[english,shorthands=off]{babel}
\usepackage{fancyhdr}
\usepackage[bookmarks=false,colorlinks=false]{hyperref}
\usepackage[protrusion=true,expansion=false]{microtype}
\emergencystretch=3em
\tolerance=600
\providecommand{\modd}{\mathrm{mod}}
\providecommand{\Disc}{\mathrm{Disc}}
\providecommand{\canont}{}
\providecommand{\asterism}{*}
\providecommand{\Hessian}{H}
\providecommand{\tome}{}
\providecommand{\page}{p.}
\pagestyle{fancy}\fancyhf{}\fancyhead[L]{Pages 501--525}\fancyhead[R]{Cayley --- Collected Papers, Vol.\ V}\renewcommand{\headrulewidth}{0.4pt}
\setlength{\headheight}{15pt}

\title{Cayley, \emph{Collected Mathematical Papers}, Vol.~V\\
PDF pp.\ 501--525\\
Papers 363 (concluded), 364, 365, 366, 367, 368, 369, 370}
\author{}
\date{}

\begin{document}
\maketitle

\noindent\textit{This installment contains the concluding pages of Paper 363 (``On the Theory of the Evolute''), and the complete texts of Paper 364 (``On a Theorem Relating to Five Points in a Plane''), Paper 365 (``On the Intersections of a Pencil of Four Lines by a Pencil of Two Lines''), Paper 366 (``Note on the Projection of the Ellipsoid''), Paper 367 (``On a Triangle In-and-Circumscribed to a Quartic Curve''), Paper 368 (``On a Problem of Geometrical Permutation''), Paper 369 (``On a Property of Commutants''), and Paper 370 (``On the Signification of an Elementary Formula of Solid Geometry'').}

\bigskip

% =====================================================================
% Paper 363 (concluded) -- printed pp. 475-479, PDF pp. 501-505
% =====================================================================
\section*{Paper 363 (concluded): On the Theory of the Evolute}

\medskip\noindent
[\textit{From the Quarterly Journal of Pure and Applied Mathematics, vol.~\textsc{vii}.\ (1866), pp.\ 169--175.}]

\medskip
I consider the particular case where the conic touches the absolute. There is no loss of generality in assuming that the contact takes place at the point $(y=0,\, z=0)$, the common tangent being therefore $z=0$; the conditions for this are $a=0$, $h=0$, and we have thence $C=0$, $F=0$. Substituting these values, the equation contains the factor $\theta$; and, throwing this out, it is
\[
X\bigl(-H\theta^{2} + (B+2G)\theta\bigr)
+ Y\bigl(A\theta^{2} - 2H\theta\bigr)
+ Z\bigl(-A\theta^{2} + 3H\theta - (B+2G)\bigr) = 0,
\]
or, what is the same thing,
\[
\theta^{2}\bigl(-HX + AY\bigr)
+ \theta\bigl((B+2G)X - 2HY - AZ\bigr)
+ \theta\bigl(3HZ\bigr)
+ \bigl(-(B+2G)Z\bigr) = 0,
\]
where it will be observed that the constant term and the coefficient of $\theta^{2}$ have the same variable factor $Z$, where $Z=0$ is the equation of the common tangent of the conic and the absolute. The evolute is in this case of the class~3. It at once appears that the line $Z=0$ is a stationary tangent of the evolute, the point of contact (or inflexion on the evolute) being given by the equations $Z=0$, $(B+2G)X - 2HY = 0$. The equation of the evolute is found by equating to zero the discriminant of the cubic function; the equation so obtained has the factor $Z$, and throwing this out the order is $=3$. The evolute is thus a curve of the class~3 and order~3, the reduction in the order from $3.2$, $=6$, to~3 being caused by the existence of an inflexion. Comparing with the former case, we see that the effect of the contact of the conic with the absolute is to give rise to an inflexion of the evolute, and to cause a reduction $=1$ in the class, and a reduction $=3$ in the order.

I return now to the general case of a curve
\[
U = (* \mid x,y,z)^{m} = 0;
\]
or rather, for greater simplicity, the equation $x^{2}+y^{2}+z^{2}=0$ for the absolute, the equation
\[
\Omega \;=\; \begin{vmatrix} X, & Y, & Z \\ x, & y, & z \\ d_{x}U, & d_{y}U, & d_{z}U \end{vmatrix} \;=\; 0;
\]
we may at once find the class of the evolute; in fact, treating $(X,Y,Z)$ as the coordinates of a given point, the two equations $U=0$, $\Omega=0$ determine the values $(x,y,z)$ of the coordinates of a point such that the normal thereof passes through

\hfill 60--2

\bigskip

\noindent the point $(X,Y,Z)$; the number of such points is the number of normals which can be drawn through a given point $(X,Y,Z)$, viz.\ it is equal to the class of the evolute. The points in question are given as the intersections of the two curves $U=0$, $\Omega=0$, which are respectively curves of the order $m$, hence the number of intersections is $=m^{2}$. It is to be observed, however, that if the curve $U=0$ has nodes or cusps, then the curve $\Omega=0$ passes through each node of the curve $U=0$, and through each cusp, the two curves having at the cusp a common tangent; that is, each node reckons for two intersections, and each cusp for three intersections. Hence, if the curve $U=0$ has $\delta$ nodes and $\kappa$ cusps, the number of the remaining points of intersection is $=m^{2}-2\delta-3\kappa$. The class of the evolute is thus $=m^{2}-2\delta-3\kappa$. The number of inflexions is in general $=0$. If, however, the given curve touches the absolute, then it has been seen in a particular case that the effect is to diminish the class by~1, and to give rise to an inflexion, the stationary tangent being in fact the common tangent of the curve and the absolute: I assume that this is the case generally. Suppose that there are $\theta$ contacts, then there will be a diminution $=\theta$ in the class, or this will be $=m^{2}-2\delta-3\kappa-\theta$; and there will be $\theta$ inflexions; there may however be special circumstances giving rise to fresh inflexions, and I will therefore assume that the number of inflexions is $=\iota'$.

Suppose in general that for any curve we have
\begin{align*}
m, &\text{ the order,}\\
n, &\text{ ,, class,}\\
\delta, &\text{ ,, number of nodes,}\\
\kappa, &\text{ ,, ,, cusps,}\\
\tau, &\text{ ,, ,, double tangents,}\\
\iota, &\text{ ,, ,, inflexions.}
\end{align*}
Then Pl\"ucker's equations give
\[
\tau - \delta = \tfrac{1}{2}(n-m)(n+m-9);
\]
\[
\iota - \kappa = 3(n-m),
\]
and we thence have
\[
\iota - \kappa + \tau - \delta = \tfrac{1}{2}(n-1)(n-2) - \tfrac{1}{2}(m-1)(m-2),
\]
or, what is the same thing,
\[
\tfrac{1}{2}(m-1)(m-2) - \delta - \kappa = \tfrac{1}{2}(n-1)(n-2) - \tau - \iota.
\]

Now M.~Clebsch in his recent paper ``Ueber die Singularit\"aten algebraischer Curven,'' \emph{Crelle}, vol.~\textsc{lxiv}.\ (1864), pp.\ 98--100, has remarked (as a consequence of the investigations of Riemann in the Integral Calculus) that whenever from a given curve another curve is derived in such manner that to each point (or tangent) of the given curve there corresponds a single tangent (or point) of the derived curve, then the expression
\[
\tfrac{1}{2}(m-1)(m-2) - \delta - \kappa \;=\; \tfrac{1}{2}(n-1)(n-2) - \tau - \iota
\]

\newpage

\noindent has the same value in the two curves respectively, or that, writing $m'$, $n'$, $\delta'$, $\kappa'$, $\tau'$, $\iota'$ for the corresponding quantities in the second curve, then we have
\[
\tfrac{1}{2}(m-1)(m-2) - \delta - \kappa \;=\; \tfrac{1}{2}(n-1)(n-2) - \tau - \iota,
\]
\[
=\tfrac{1}{2}(m'-1)(m'-2) - \delta' - \kappa' \;=\; \tfrac{1}{2}(n'-1)(n'-2) - \tau' - \iota';
\]
and consequently that, knowing any two of the quantities $m'$, $n'$, $\delta'$, $\kappa'$, $\tau'$, $\iota'$, the remainder of them can be determined by means of this relation and of Pl\"ucker's equations. The theorem is applicable to the evolute according to the foregoing generalized definition${}^{(1)}$; and starting from the values
\[
n' = m^{2} - 2\delta - 3\kappa - \theta,
\]
we find in the first instance
\[
\iota' = \tfrac{1}{2}(n'-1)(n'-2) - \tfrac{1}{2}(m-1)(m-2) + \delta + \kappa - \iota;
\]
and substituting in the equation
\[
m' = n'(n'-1) - 2\tau' - 3\iota',
\]
we find
\[
m' = 2(n'-1) + (m-1)(m-2) - 2\delta - 2\kappa - \iota;
\]
and the equation $\iota' - \kappa' = 3(n'-m')$ gives also
\[
\kappa' = -3(n'-m') + \iota',
\]
whence, attending to the value of $n'$, we find the following system of equations for the singularities of the evolute, viz.
\begin{align*}
n' &= m^{2} - 2\delta - 3\kappa - \theta,\\
m' &= 3m(m-1) - 6\delta - 8\kappa - 2\theta - \iota,\\
\iota' &= \iota,\\
\kappa' &= 3m(2m-3) - 12\delta - 15\kappa - 3\theta - 2\iota,
\end{align*}
and the values of $\tau'$ and $\delta'$ may then also be found from the equations
\[
m = n(n-1) - 2\tau' - 3\iota',
\]
\[
n' = m'(m'-1) - 2\delta' - 3\kappa'.
\]

I have given the system in the foregoing form, as better exhibiting the effect of the inflexions; but as each of the $\theta$ contacts with the absolute gives an inflexion, we

\medskip
\footnoterule
\noindent\footnotesize ${}^{1}$ M.~Clebsch in fact applies it to the evolute in the ordinary sense of the term, but by inadvertently assuming $\iota' = \kappa$ instead of $\iota' = \theta$ he is led to some incorrect results.

\newpage

\noindent\normalsize may write $\iota = \theta + \iota''$, where, in the absence of special circumstances giving rise to any more inflexions, $\iota''=0$. The system thus becomes
\begin{align*}
n' &= m^{2} - 2\delta - 3\kappa - \theta,\\
m' &= 3m(m-1) - 6\delta - 8\kappa - 3\theta - \iota'',\\
\iota' &= \theta + \iota'',\\
\kappa' &= 3m(2m-3) - 12\delta - 15\kappa - 5\theta - 2\iota'',
\end{align*}
so that each contact with the absolute diminishes the class by~1, the order by~3, and the number of cusps by~5.

I remark that when the absolute becomes a pair of points, a contact of the given curve $m$ means one of two things: either the curve touches the line through the two points, or else it passes through one of the two points: the effect of a contact of either kind is as above stated. Suppose that the two points are the circular points at infinity, and let $m=2$, the evolute in question is then the evolute of a conic, in the ordinary sense of the word evolute. We have, in general, class $=4$, order $=6$; but if the conic touches the line infinity (that is, in the case of the parabola), the reductions are 1 and 3, and we have class $=3$, order $=3$, which is right. If the conic passes through one of the circular points of infinity, then in like manner the reductions are 1 and 3; and therefore if the conic passes through each of the circular points at infinity (that is, in the case of a circle), the reductions are 2 and 6, and we have class $=2$, order $=0$, which is also right; for the evolute is in this case the centre, regarded as a pair of coincident points. That this is so, or that the class is to be taken to be (not $=1$ but) $=2$, appears by the consideration that the number of normals to the circle from a given point is in fact $=2$, the two normals being, however, coincident in position.

To complete the theory in the general case where the absolute is a proper conic, I remark that, besides the inflexions which arise from contacts of the given curve with the absolute, there will be an inflexion, first, for each stationary tangent of the given curve which is also a tangent of the absolute; secondly, for each cusp of the given curve situate on the absolute. Hence, if the number of such stationary tangents be $=\lambda$, and the number of such cusps be $=\mu$, we may write $\iota'' = \lambda + \mu$, and therefore also $\iota = \theta + \lambda + \mu$.

I remark also that we have
\[
-2\delta - 8\kappa = -m(m-1) + n,
\]
and therefore also
\[
-6\delta - 8\kappa = -3m(m-2) + \iota,
\]
\[
-12\delta - 15\kappa = -6m^{2} + 15m - 3n + 3\iota.
\]
The general formul\ae{} thus become
\begin{align*}
n' &= m + n - \theta,\\
m' &= 3n + \iota - 2\theta - \iota'',\\
\iota' &= \iota,\\
\kappa' &= 6n - 3m + 3\iota - 3\theta - 2\iota''.
\end{align*}

\newpage

\noindent If instead of the given curve we consider its reciprocal in regard to the absolute, then
\[
m,\, n,\, \delta,\, \kappa,\, \tau,\, \iota;\; \theta,\, \lambda,\, \mu;\; \iota,\, = \theta + \lambda + \mu
\]
are changed into
\[
n,\, m,\, \tau,\, \iota,\, \delta,\, \kappa;\; \theta,\, \mu,\, \lambda;\; \iota,\, = \theta + \mu + \lambda
\]
respectively.

Hence for the evolute of the reciprocal curve we have
\begin{align*}
n' &= n + m - \theta,\\
m' &= 3n + \kappa - 2\theta - \iota'',\\
\iota' &= \iota,\\
\kappa' &= 6n - 3m + 3\kappa - 3\theta - \iota'',
\end{align*}
which, attending to the relation $\iota - \kappa = 3(n - m)$, are in fact the same as the former values; that is, the evolute of the given curve, and the evolute of the reciprocal curve are curves of the same class and order, and which have the same singularities.

\medskip
\noindent\textit{Cambridge, February 22, 1865.}

\bigskip

% =====================================================================
% Paper 364 -- printed pp. 480-483, PDF pp. 506-509
% =====================================================================
\section*{Paper 364: On a Theorem Relating to Five Points in a Plane}

\medskip\noindent
[\textit{From the Philosophical Magazine, vol.\ \textsc{xxix}.\ (1865), pp.\ 460--464.}]

\medskip
Two triangles, $ABC$, $A'B'C'$ which are such that the lines $AA'$, $BB'$, $CC'$ meet in a point, are said to be in perspective; and a triangle $A'B'C'$, the angles $A'$, $B'$, $C'$ of which lie in the sides $BC$, $CA$, $AB$ respectively, is said to be inscribed in the triangle $ABC$; hence, if $A'$, $B'$, $C'$ are the intersections of the sides by the lines $AO$, $BO$, $CO$ respectively (where $O$ is any point whatever), the triangle $A'B'C'$ is said to be perspectively inscribed in the triangle $ABC$, viz.\ it is so inscribed by means of the point $O$.

We have the following theorem, relating to any triangle $ABC$, and two points $O$, $O'$. If in the triangle $ABC$, by means of the point $O$, we inscribe a triangle $A'B'C'$, and in the triangle $A'B'C'$, by means of the point $O'$, we inscribe a triangle $\alpha\beta\gamma$, then the triangles $ABC$, $\alpha\beta\gamma$ are in perspective, viz.\ the lines $A\alpha$, $B\beta$, $C\gamma$ will meet in a point.

This is very easily proved analytically; in fact, taking $x=0$, $y=0$, $z=0$ for the equations of the lines $B'C'$, $C'A'$, $A'B'$ respectively, and $(X, Y, Z)$ for the coordinates of the point $O$, then the coordinates of $(A, B, C)$ are found to be $(-X, Y, Z)$, $(X, -Y, Z)$, $(X, Y, -Z)$ respectively. Moreover, if $(X', Y', Z')$ are the coordinates of the point $O'$, then the coordinates of $(\alpha, \beta, \gamma)$ are found to be $(0, Y', Z')$, $(X', 0, Z')$, $(X', Y', 0)$ respectively. Hence the equations of the lines $A\alpha$, $B\beta$, $C\gamma$ are
\[
\begin{vmatrix} x, & y, & z \\ -X, & Y, & Z \\ 0, & Y', & Z' \end{vmatrix} = 0,\quad
\begin{vmatrix} x, & y, & z \\ X, & -Y, & Z \\ X', & 0, & Z' \end{vmatrix} = 0,\quad
\begin{vmatrix} x, & y, & z \\ X, & Y, & -Z \\ X', & Y', & 0 \end{vmatrix} = 0,
\]

\newpage

\noindent that is
\begin{align*}
x(YZ' - Y'Z) + y(Z'X) + z(-XY') &= 0,\\
x(-YZ') + y(ZX' - Z'X) + z(X'Y) &= 0,\\
x(+Y'Z) + y(-ZX') + z(XY' - X'Y) &= 0,
\end{align*}
which are obviously the equations of three lines which meet in a point.

But the theorem may be exhibited as a theorem relating to a quadrangle $1234$ and a point $O'$; for writing $1$, $2$, $3$, $4$ in place of $A$, $B$, $C$, $O$, the triangle $A'B'C'$ is in fact the triangle formed by the three centres $41.23$, $42.31$, $43.12$ of the quadrangle $1234$, hence the triangle in question must be similarly related to each of the four triangles $423$, $431$, $412$, $123$; or, forming the diagram
\[
\begin{array}{c|cccc}
 & P & Q & R & S \\\hline
41.23 & 4 & 3 & 2 & 1 \\
42.31 & 3 & 4 & 1 & 2 \\
43.12 & 2 & 1 & 4 & 3
\end{array}
\]
we have the following form of the theorem: viz.\ the lines
\begin{align*}
\alpha 4,\; \beta 3,\; \gamma 2 \text{ meet in a point } P,&\\
\alpha 3,\; \beta 4,\; \gamma 1 \;\;,,\;\; ,, \;\; Q,&\\
\alpha 2,\; \beta 1,\; \gamma 4 \;\;,,\;\; ,, \;\; R,&\\
\alpha 1,\; \beta 2,\; \gamma 3 \;\;,,\;\; ,, \;\; S,&
\end{align*}
or, what is the same thing, we have with the points $1$, $2$, $3$, $4$ and the point $O'$ constructed the four points $P$, $Q$, $R$, $S$ such that
\begin{align*}
1S,\; 2R,\; 3Q,\; 4P \text{ meet in a point } \alpha,&\\
2S,\; 1R,\; 4Q,\; 3P \;\;,,\;\; ,, \;\; \beta,&\\
3S,\; 4R,\; 1Q,\; 2P \;\;,,\;\; ,, \;\; \gamma.&
\end{align*}

The eight points $1$, $2$, $3$, $4$, $P$, $Q$, $R$, $S$ form a figure such as the perspective representation of a parallelopiped, or, if we please, a cube; and not only so, but the

\medskip
\textit{[Figure: schematic of a cube in perspective, with vertices labelled $1,2,3,4$ on the lower face and $P,Q,R,S$ on the upper face.]}

\medskip
plane figure is really a certain perspective representation of the cube; this identification depends on the following two theorems:

\hfill 61

\newpage

\noindent 1. Considering the four summits $1$, $2$, $3$, $4$, which are such that no two of them belong to the same edge, then, if through any point $O$ we draw
\begin{align*}
& \text{the line } OA' \text{ meeting the lines } 41,\, 23,\\
& \;\; ,, \;\; OB' \;\; ,, \;\; ,, \;\; 42,\, 31,\\
& \;\; ,, \;\; OC' \;\; ,, \;\; ,, \;\; 43,\, 12,
\end{align*}
and the lines $O\alpha$, $O\beta$, $O\gamma$ parallel to the three edges of the cube respectively, the three planes $(OA', O\alpha)$, $(OB', O\beta)$, $(OC', O\gamma)$ will meet in a line.

2. For a properly selected position of the point $O$,
\begin{align*}
& \text{the lines } OB',\; OC',\; O\alpha \text{ will lie in a plane,}\\
& \;\; ,, \;\; OC',\; OA',\; O\beta \;\; ,,\\
& \;\; ,, \;\; OA',\; OB',\; O\gamma \;\; ,,
\end{align*}

In fact for such a position of $O$, projecting the whole figure on any plane whatever, the lines $O1$, $O2$, $O3$, $O4$, $OP$, $OQ$, $OR$, $OS$, $O\alpha$, $O\beta$, $O\gamma$, $OA'$, $OB'$, $OC'$ meet the plane of projection in the points $1$, $2$, $3$, $4$, $P$, $Q$, $R$, $S$, $\alpha$, $\beta$, $\gamma$, $A'$, $B'$, $C'$ related to each other as in the last-mentioned form of the plane theorem. To prove the two solid theorems, take $O$ for the origin, $O\alpha$, $O\beta$, $O\gamma$ for the axes, $(\alpha, \beta, \gamma)$ for the coordinates of the summit $S$, and 1 for the edge of the cube,
\begin{align*}
& \text{the coordinates of } 1 \text{ are } \alpha+1,\, \beta,\, \gamma,\\
& \;\;\;\;\;\;\;\;\;\;\;\;\;\;\;\;\;\;\;\;\;\;\;\; 2 \;\;,, \;\; \alpha,\, \beta+1,\, \gamma,\\
& \;\;\;\;\;\;\;\;\;\;\;\;\;\;\;\;\;\;\;\;\;\;\;\; 3 \;\;,, \;\; \alpha,\, \beta,\, \gamma+1,\\
& \;\;\;\;\;\;\;\;\;\;\;\;\;\;\;\;\;\;\;\;\;\;\;\; 4 \;\;,, \;\; \alpha+1,\, \beta+1,\, \gamma+1.
\end{align*}
The equations of the line $OA'$, or say of the line $O(41, 23)$, are those of the planes $O41$, $O23$, viz.\ they are
\[
\begin{vmatrix} x, & y, & z \\ \alpha+1, & \beta, & \gamma \\ \alpha+1, & \beta+1, & \gamma+1 \end{vmatrix} = 0,\qquad
\begin{vmatrix} x, & y, & z \\ \alpha, & \beta+1, & \gamma \\ \alpha, & \beta, & \gamma+1 \end{vmatrix} = 0;
\]
that is
\[
x(\beta - \gamma) - (\alpha + 1)(y - z) = 0,
\]
\[
x(\beta + \gamma + 1) - \alpha(y + z) = 0.
\]
Writing for shortness
\[
M = \alpha + \beta + \gamma + 1,
\]
these equations give
\[
x : y : z \;=\; \frac{2\alpha(\alpha + 1)}{(M + 2\gamma\alpha)(M + 2\alpha\beta)} : \frac{1}{M + 2\gamma\alpha} : \frac{1}{M + 2\alpha\beta},
\]
or, completing the system,
\[
\text{for line } OA' \text{ we have } \; x : y : z \;=\; \frac{2\alpha(\alpha + 1)}{(M + 2\gamma\alpha)(M + 2\alpha\beta)} : \frac{1}{M + 2\gamma\alpha} : \frac{1}{M + 2\alpha\beta};
\]

\newpage

\noindent for line $OB'$ we have
\[
x : y : z \;=\; \frac{1}{M + 2\beta\gamma} : \frac{2\beta(\beta + 1)}{(M + 2\alpha\beta)(M + 2\beta\gamma)} : \frac{1}{M + 2\alpha\beta};
\]
for line $OC'$ we have
\[
x : y : z \;=\; \frac{1}{M + 2\beta\gamma} : \frac{1}{M + 2\gamma\alpha} : \frac{2\gamma(\gamma + 1)}{(M + 2\beta\gamma)(M + 2\gamma\alpha)}.
\]

The equations of the lines $O\alpha$, $O\beta$, $O\gamma$ are of course $(y=0, z=0)$, $(z=0, x=0)$, $(x=0, y=0)$ respectively; and we therefore see at once that the planes $(OA', O\alpha)$, $(OB', O\beta)$, $(OC', O\gamma)$ meet in a line, viz.\ in the line which has for its equations
\[
\frac{x}{M + 2\beta\gamma} \;=\; \frac{y}{M + 2\gamma\alpha} \;=\; \frac{z}{M + 2\alpha\beta}.
\]

The lines $OB'$, $OC'$, $O\alpha$ will lie in a plane, if only
\[
\frac{1}{(M + 2\beta\gamma)^{2}} \;=\; \frac{4\beta\gamma(\beta + 1)(\gamma + 1)}{(M + 2\alpha\beta)(M + 2\gamma\alpha)(M + 2\beta\gamma)^{2}},
\]
that is
\[
(M + 2\beta\gamma)^{2} \;=\; 4\beta\gamma(\beta + 1)(\gamma + 1),
\]
or, as this may be written,
\[
M^{2} + 4\beta\gamma(\alpha + \beta + \gamma + 1 + \beta\gamma) \;=\; 4\beta\gamma(\beta\gamma + \beta + \gamma + 1);
\]
that is
\[
M^{2} + 4\alpha\beta\gamma \;=\; 0,
\]
\[
(\alpha + \beta + \gamma + 1)^{2} + 4\alpha\beta\gamma \;=\; 0;
\]
and from the symmetry of this equation we see that, when it is satisfied,
\begin{align*}
& \text{the lines } OB',\; OC',\; O\alpha \text{ will lie in a plane,}\\
& \;\; ,, \;\; OC',\; OA',\; O\beta \;\; ,,\\
& \;\; ,, \;\; OA',\; OB',\; O\gamma \;\; ,,;
\end{align*}
viz.\ this will be the case when the point $O$ is situate in the cubic surface represented by the last-mentioned equation; this completes the demonstration of the solid theorems.

It is clear that considering five points $1$, $2$, $3$, $4$, $5$ in a plane, then, since any one of these may be taken for the point $O'$ of the foregoing theorem, the theorem exhibited in the first instance as a theorem relating to a triangle and two points, and afterwards as a theorem relating to a quadrangle and a point, is really a theorem relating to five points in a plane. There are, of course, five different systems of points $(P, Q, R, S)$, corresponding to the different combinations of four out of the five points.

\medskip
\noindent\textit{Cambridge, March 6, 1865.}

\hfill 61--2

\newpage

% =====================================================================
% Paper 365 -- printed pp. 484-486, PDF pp. 510-512
% =====================================================================
\section*{Paper 365: On the Intersections of a Pencil of Four Lines by a Pencil of Two Lines}

\medskip\noindent
[\textit{From the Philosophical Magazine, vol.\ \textsc{xxix}.\ (1865), pp.\ 501--503.}]

\medskip
Pl\"ucker has considered (``Analytisch-geometrische Aphorismen,'' \emph{Crelle}, vol.\ \textsc{xi}.\ (1834) pp.\ 26--32) the theory of the eight points which are the intersections of a pencil of four lines by any two lines, or say the intersections of a pencil of four lines by a pencil of \emph{two} lines: viz., the eight points may be connected two together by twelve new lines; the twelve lines meet two together in forty-two new points; and of these, six lie on a line through the centre of the two-line pencil, twelve lie four together on three lines through the centre of the four-line pencil, and twenty-four lie two together on twelve lines, also through the centre of the four-line pencil.

The first and third of these theorems, viz.\ (1) that the six points lie on a line through the centre of the two-line pencil, and (3) that the twenty-four points lie two together on twelve lines through the centre of the four-line pencil, belong to the more simple theory of the intersections of a pencil of three lines by a pencil of \emph{two} lines; the second theorem, viz.\ (2) the twelve points lie four together on three lines through the centre of the four-line pencil, is the only one which properly belongs to the theory of the intersections of a pencil of four lines by a pencil of \emph{two} lines. The theorem in question (proved analytically by Pl\"ucker) may be proved geometrically by means of two fundamental theorems of the geometry of position: these are the theorem of two triangles in perspective, and Pascal's theorem for a line-pair. I proceed to show how this is.

Consider a pencil of two lines meeting a pencil of four lines in the eight points $(a, b, c, d)$, $(a', b', c', d')$; so that the two lines are $abcd$, $a'b'c'd'$, meeting suppose in $Q$; and the four lines are $aa'$, $bb'$, $cc'$, $dd'$, meeting suppose in $P$; then the twelve points are
\begin{align*}
a'd \cdot c'b,\; ad' \cdot cb',\; a'c \cdot d'b,\; ac' \cdot db'\;\; & \text{lying in a line through } P,\\
a'b \cdot d'c,\; ab' \cdot dc',\; a'd \cdot b'c,\; ad' \cdot bc'\;\; & \text{,, ,, ,,}\\
a'c \cdot b'd,\; ac' \cdot bd',\; a'b \cdot c'd,\; ab' \cdot cd'\;\; & \text{,, ,, ,,;}
\end{align*}
where the combinations are most easily formed as follows; viz., for the first four points starting from the arrangement $\dfrac{a\;c}{d\;b}$ (or any other arrangement having the diagonals $ab \cdot cd$), and thence writing down the four expressions
\[
\begin{array}{l}
a'c,\; ac',\; a'c,\; ac'\\
d'b,\; d'b',\; d'b,\; db',
\end{array}
\]
we read off from these the symbols of the four points; and the like for the other two sets of four points.

Now, considering the points $(a, b, c)$ and $(a', b', c')$, the points $ab' \cdot a'b$, $ac' \cdot a'c$, $bc' \cdot b'c$ lie in a line through $Q$; and similarly the points $ab' \cdot a'b$, $ad' \cdot a'd$, $bd' \cdot b'd$ lie in a line through $Q$; which lines, inasmuch as they each contain the points $Q$ and $ab' \cdot a'b$, must be one and the same line; considering the combinations $(b, c, d)$, $(b', c', d')$, the line in question also passes through $cd' \cdot c'd$; that is, the six points $ab' \cdot a'b$, $ac' \cdot a'c$, $ad' \cdot a'd$, $bc' \cdot b'c$, $bd' \cdot b'd$, $cd' \cdot c'd$ lie in a line through $Q$, which is in fact the before-mentioned first theorem. Hence the points $ab' \cdot a'b$ and $cd' \cdot c'd$ lie in a line through $Q$; or, calling these points $M$ and $N$ respectively, the triangles $Maa'$, $Mbb'$, $Ncc'$, $Ndd'$ are in perspective. Hence, considering the two triangles $Maa'$, $Ndd'$ (or, if we please, the complementary set $Mbb'$, $Ncc'$), the corresponding sides are
\begin{align*}
Ma\;,\; Nd \;&\text{meeting in } ab' \cdot dc',\\
Ma',\; Nd' \;&\;\; ,, \;\; a'b \cdot d'c,\\
aa'\;,\; dd' \;&\;\; ,, \;\; P;
\end{align*}
that is, the points $ab' \cdot dc'$, $a'b \cdot d'c$ lie in a line through $P$.

Similarly $ad' \cdot a'd$ and $bc' \cdot b'c$ lie in a line through $Q$; or, calling these points $H$, $I$ respectively, the triangles $Haa'$, $Hdd'$, $Ibb'$, $Icc'$ are in perspective; and considering the combination $Hdd'$, $Ibb'$ (or, if we please, the complementary set $Haa'$, $Icc'$), the corresponding sides are
\begin{align*}
Ha\;,\; Ib \;&\text{meeting in } ad' \cdot bc',\\
Ha',\; Ib' \;&\;\; ,, \;\; a'd \cdot c'b,\\
aa'\;,\; bb' \;&\;\; ,, \;\; P;
\end{align*}
that is, the points $a'd \cdot c'b$, $ad' \cdot cb'$ lie in a line through $P$.

\newpage

\noindent It remains to be shown that the two lines through $P$, viz.\ the line containing $ab' \cdot dc'$ and $a'b \cdot d'c$, and the line containing $ad' \cdot bc'$ and $a'd \cdot cb'$, are one and the same line. This will be the case if, for instance, $ab' \cdot dc'$ and $ad' \cdot bc'$ also lie in a line through $P$.

We have the points $(a, b, d)$ in a line, and the points $(b', c', d')$ in a line; the points $a$, $d$, $b'$, $c'$ are also called $A$, $B'$, $B$, $A'$ respectively; $ad'$, $bb'$ meet in $C$, and $bc'$, $dd'$ meet in $C'$; hence, considering the hexagon $ad'db'bc'$, the lines
\begin{align*}
ad',\; b'b \;&\text{meet in } C,\\
d'd,\; bc' \;&\;\; ,, \;\; C',\\
db',\; ca' \;&\;\; ,, \;\; AA' \cdot BB';
\end{align*}
and hence these three points lie in a line; or, what is the same thing, the lines $AA'$, $BB'$, and $CC'$ meet in a point; that is, the triangles $ABC$, $A'B'C'$ are in perspective: the corresponding sides are
\begin{align*}
AB,\; A'B'\;&\text{ that is, } ab',\; c'd, \text{ meeting in } ab' \cdot c'd,\\
BC,\; B'C'\;&\;\; ,, \;\; b'b,\; d'd,\;\; ,, \;\; P,\\
CA,\; C'A'\;&\;\; ,, \;\; ad',\; bc',\;\; ,, \;\; ad' \cdot bc';
\end{align*}
and these three points lie in a line; that is, the points $ab' \cdot dc'$ and $ad' \cdot bc'$ lie in a line through $P$. Hence the line through $ab' \cdot dc'$ and $a'b \cdot d'c$ and the line through $ad' \cdot bc'$ and $a'd \cdot cb'$ are one and the same line; that is,
\[
\text{the points } ab' \cdot dc',\; a'b \cdot d'c,\; ad' \cdot bc',\; a'd \cdot b'c \text{ lie in a line through } P.
\]
This proves the existence of one of the lines through $P$; and that of the other two lines follows from the symmetry of the figure; it thus appears that the twelve points lie four together on three lines through $P$.

\medskip
\noindent\textit{Cambridge, April 11, 1865.}

\bigskip

% =====================================================================
% Paper 366 -- printed pp. 487-488, PDF pp. 513-514
% =====================================================================
\section*{Paper 366: Note on the Projection of the Ellipsoid}

\medskip\noindent
[\textit{From the Philosophical Magazine, vol.\ \textsc{xxx}.\ (1865), pp.\ 50--52.}]

\medskip
Consider an ellipsoid, situate any way whatever in regard to the eye and the plane of the picture; the apparent contour of the ellipsoid is an ellipse, the intersection of the plane of the picture by the tangent cone having the eye for vertex; this cone touches the ellipsoid along a plane curve (the intersection of the ellipsoid by the polar plane of the eye), which may be called the contour section; and the apparent contour is thus the projection of the contour section. Consider any other plane section; the projection thereof has double contact (real or imaginary) with the projection of the contour section: the common tangents are the intersections with the plane of the picture of the tangent planes of the tangent cone which pass through the pole of the section; or, what is the same thing, they are the tangents to the projection of the contour section, or to the projection of the section, from the point which is the projection of the pole of the section. The projection of the pole lies in the line which is the projection of the diameter conjugate to the plane of the section; and in particular, if the section is central, that is, if the plane thereof passes through the centre of the ellipsoid, then the pole is the point at infinity on the conjugate diameter; whence also if the eye be at an infinite distance, so that the projection is a projection by parallel rays, then the projection of the pole is the point at infinity on the projection of the conjugate diameter; and therefore the common tangents of the projections of the section and the contour section are in this case parallel to the projection of the diameter conjugate to the plane of the section.

Suppose that the plane of the picture is parallel to a principal plane of the ellipsoid, and that the projection is by parallel rays; then if $OA$, $OB$, $OC$ are the projections of the semiaxes ($OA$, $OC$ will be at right angles to each other if the plane parallel to the plane of the picture is that of $xz$), the projections of the principal sections are the ellipses having for conjugate semidiameters $OB$, $OC$; $OC$, $OA$; $OA$, $OB$ respectively. Hence to the ellipse $OB$, $OC$ drawing the two tangents which are parallel to $OA$, to the ellipse $OC$, $OA$ the two tangents which are parallel to $OB$, and to the ellipse $OA$, $OB$ the two tangents which are parallel to $OC$, we have on each of these ellipses the two points which are the points of contact therewith of the ellipse which is the projection of the contour section, or apparent contour of the ellipsoid; that is, we know six points, and at each of these points the tangent, of the last-mentioned ellipse; and the ellipse in question, or apparent contour of the ellipsoid, can thus be traced by hand accurately enough for ordinary purposes.

In connexion with what precedes, I may notice a convenient construction for the projection of a circle. Suppose that we have given the projection of the circumscribed square $ABCD$; then if we know the projection of one of the points $M$, $N$, $P$, $Q$, say of the point $M$, the projections of all the points and lines of the figure can be obtained graphically by the ruler only with the utmost facility; that is, in the ellipse which is the projection of the circle we have eight points, and the tangent at each of them, and the ellipse may then be drawn by hand. And to find the projection of the point $M$, it is only necessary to remark that in the figure the anharmonic ratio $\dfrac{AM}{MO} : \dfrac{AC}{OC}$ of the points $A$, $M$, $O$, $C$ is $=\tfrac{1}{2}(\sqrt{2}-1)$; hence the corresponding anharmonic ratio of the projections of the four points is also $=\tfrac{1}{2}(\sqrt{2}-1)$; and the projections of $A$, $B$, $C$, $D$, and consequently those of $A$, $C$, $O$, being known, the projection of $M$ is thus also known.

\medskip
\textit{[Figure: a square $ABCD$ inscribing a circle, with midpoints $M$, $N$, $P$, $Q$ of the sides, the centre $O$, and diagonals through $O$.]}

\medskip
\noindent\textit{Cambridge, June 15, 1865.}

\bigskip

% =====================================================================
% Paper 367 -- printed pp. 489-492, PDF pp. 515-518
% =====================================================================
\section*{Paper 367: On a Triangle In-and-Circumscribed to a Quartic Curve}

\medskip\noindent
[\textit{From the Philosophical Magazine, vol.\ \textsc{xxx}.\ (1865), pp.\ 340--342.}]

\medskip
The quartic curve $(x^{2} - a^{2})^{2} + (y^{2} - b^{2})^{2} = c^{4}$ presents a simple example of a triangle in-and-circumscribed to a single curve, viz.\ such that each angle of the triangle is situate on, and each side touches, the curve. Assuming that the triangle is symmetrically situate in regard to the axis of $y$, viz.\ if it be the isosceles triangle $ace$, the sides whereof touch the curve in the points $B$, $D$, $F$ respectively, then we must have a single relation between the constants $a$, $b$, $c$ of the curve; or if (as may be done without loss of generality) we write $a=1$, then there must be a single relation between $b$ and $c$. The relation in question is most conveniently expressed by putting $b$ and $c$ equal to certain functions of a parameter $\phi$, which is in fact $=\tan^{2}\theta$, if $\theta$ be the angle at the base of the triangle; the equation of the curve is thus obtained in the form
\[
(x^{2} - 1)^{2} + \left( y^{2} - \frac{\phi^{4} + 4\phi^{2} - 1}{4\phi(\phi^{2} + 1)} \right)^{2} \;=\; 1 + \frac{(\phi^{2} - 1)^{4}}{16\phi^{2}(\phi^{2} + 1)^{2}};
\]

\medskip
\textit{[Figure: an isosceles triangle $ace$ with apex $a$ and base $ce$, with the points $B$, $D$, $F$ marked on the sides $ac$, $ce$, $ea$ respectively.]}

\hfill 62

\newpage

\noindent and the coordinates of $a$, $c$, $e$, $B$, $D$, $F$ are as follows:
\begin{align*}
& \text{Coordinates of } a \text{ are } 0,\; \sqrt{\tfrac{1}{2}\phi},\\
& \;\;\;\;\;\;\;\;\;\;\;\;\;\;\;\; c,e, \;\;,, \;\; \pm \sqrt{2},\; -\sqrt{\tfrac{1}{2}\phi},\\
& \;\;\;\;\;\;\;\;\;\;\;\;\;\;\;\; D,\;\;,, \;\; 0,\; -\sqrt{\tfrac{1}{2}\phi},\\
& \;\;\;\;\;\;\;\;\;\;\;\;\;\;\;\; B,F \;\;,, \;\; \pm \frac{\phi^{2} - 1}{\sqrt{2}(\phi^{2} + 1)},\; \frac{2\sqrt{\phi}}{\sqrt{2}(\phi^{2} + 1)}.
\end{align*}

It is easy to verify that the points $a$, $c$, $e$, $D$ are points of the curve, and it is obvious that the tangent at $D$ is the horizontal line $ce$. It only remains to be shown that $B$ and $F$ are points of the curve, and that the tangents at these points are the lines $ac$ and $ea$ respectively. It is sufficient to consider one of the two points, say the point $F$; and taking its coordinates to be
\[
\xi = \frac{\phi^{2} - 1}{\sqrt{2}(\phi^{2} + 1)},\qquad \eta = \frac{2\sqrt{\phi}}{\sqrt{2}(\phi^{2} + 1)},
\]
we have to show that $(\xi, \eta)$ is a point of the curve, and that the equation of the tangent at this point is $X\sqrt{\phi} + Y = \sqrt{\tfrac{1}{2}\phi}$, where $(X, Y)$ are current coordinates.

First, to show that $(\xi, \eta)$ is a point of the curve, the equation to be verified may be written
\[
(\xi^{2} - 1)^{2} + \left( \eta^{2} - \frac{\phi^{4} + 4\phi^{2} - 1}{4\phi(\phi^{2} + 1)} \right)^{2} \;=\; 1 + \frac{(\phi^{2} - 1)^{4}}{16\phi^{2}(\phi^{2} + 1)^{2}},
\]
and we have
\[
\xi^{2} - 1 = -\frac{\phi^{4} + 6\phi^{2} + 1}{2(\phi^{2} + 1)^{2}},\qquad \eta^{2} - \frac{\phi^{4} + 4\phi^{2} - 1}{4\phi(\phi^{2} + 1)} = -\frac{(\phi^{2} - 1)(\phi^{4} + 6\phi^{2} + 1)}{4\phi(\phi^{2} + 1)^{2}},
\]
so that the equation becomes
\[
\frac{(\phi^{4} + 6\phi^{2} + 1)^{2}}{4(\phi^{2} + 1)^{4}} + \frac{(\phi^{2} - 1)^{2}(\phi^{4} + 6\phi^{2} + 1)^{2}}{16\phi^{2}(\phi^{2} + 1)^{4}} \;=\; \frac{(\phi^{4} + 6\phi^{2} + 1)^{2}}{16\phi^{2}(\phi^{2} + 1)^{2}};
\]
that is
\[
4\phi^{2} + (\phi^{2} - 1)^{2} \;=\; (\phi^{2} + 1)^{2},
\]
which is right.

Next, the equation of the tangent at the point $(\xi, \eta)$ is
\[
(\xi^{2} - a^{2})(\xi X - a^{2}) + (\eta^{2} - b^{2})(\eta Y - b^{2}) - c^{4} \;=\; 0;
\]
or, substituting for $\xi$, $\eta$, $\xi^{2} - 1$, and $\eta^{2} - \dfrac{\phi^{4} + 4\phi^{2} - 1}{4\phi(\phi^{2} + 1)}$, their values, and throwing out a factor $\dfrac{\phi^{4} + 6\phi^{2} + 1}{16\phi^{2}(\phi^{2} + 1)^{3}}$, the equation becomes

\newpage

\noindent
\[
-8\phi^{2}\left( X\frac{\phi^{2} - 1}{\sqrt{2}(\phi^{2} + 1)} - 1 \right) - 4\phi(\phi^{2} - 1)\left( Y\frac{2\sqrt{\phi}}{\sqrt{2}(\phi^{2} + 1)} - \frac{\phi^{4} + 4\phi^{2} - 1}{4\phi(\phi^{2} + 1)} \right) \;=\; \phi^{4} + 6\phi^{2} + 1,
\]
or, what is the same thing,
\[
-8\phi^{2}\{ X(\phi^{2} - 1) - \sqrt{2}(\phi^{2} + 1) \} - (\phi^{2} - 1)\{ Y \cdot 8\phi\sqrt{\phi} - \sqrt{2}(\phi^{4} + 4\phi^{2} - 1) \},
\]
that is
\[
(\phi^{2} - 1)(-8\phi^{2} X - 8\phi\sqrt{\phi}\, Y)
\]
\[
= \sqrt{2}(\phi^{2} + 1)(\phi^{4} + 6\phi^{2} + 1) - \sqrt{2}(\phi^{2} + 1) \cdot 8\phi^{2} - \sqrt{2}(\phi^{2} - 1)(\phi^{4} + 4\phi^{2} - 1),
\]
\[
= \sqrt{2}(\phi^{2} + 1)(\phi^{2} - 1)^{2} - \sqrt{2}(\phi^{2} - 1)(\phi^{4} + 4\phi^{2} - 1),
\]
\[
= \sqrt{2}(\phi^{2} - 1)\{(\phi^{2} - 1)^{2} - (\phi^{4} + 4\phi^{2} - 1)\},
\]
\[
= -4\sqrt{2}(\phi^{2} - 1)\phi^{2},
\]
whence, finally,
\[
X\sqrt{\phi} + Y \;=\; \sqrt{\tfrac{1}{2}\phi},
\]
which is the required equation.

It may be remarked that for $\phi=1$, the equation of the curve is $(x^{2} - 1)^{2} + (y^{2} - \tfrac{1}{2})^{2} = 1$, which is the binodal form $a^{2} > b^{2}$, $c^{4} = a^{4}$. We have in this case $\xi = 0$, $\eta = \sqrt{\tfrac{1}{2}}$, and the curve and triangle are as shown in the figure, viz.\ the base $ce$ of the triangle, instead of being a proper tangent, is a line through the node $D$. For any other value of $\phi$, the curve consists of an exterior oval (pinched in at the sides and the top and bottom) and of an interior oval; the angles $a$, $c$, $e$ lie in the exterior oval, the sides $ac$, $ea$ touch the interior oval, and the base $ce$ touches the exterior oval.

\medskip
\textit{[Figure: an exterior oval pinched at the sides and top/bottom, containing a smaller interior oval; the triangle $ace$ is inscribed with vertices on the exterior oval and sides $ac$, $ea$ tangent to the interior oval, base $ce$ tangent to the exterior oval.]}

\hfill 62--2

\newpage

\noindent If, to fix the ideas, we assume $\phi > 1$, then we have always $c^{2} > a^{2} < a^{2} + b^{2}$: for $\phi = 1$ we have, as appears above, $b^{2} = \tfrac{1}{2}$, which is $< a^{2}$; but for a certain value of $\phi$ between 3 and 4, $b^{2}$ becomes $=a^{2}$, and for any greater value of $\phi$ we have $b^{2} > a^{2}$. The condition for the equality of $a^{2}$ and $b^{2}$ is
\[
\phi^{4} + 4\phi^{2} - 1 \;=\; 4\phi(\phi^{2} + 1),\quad\text{or}\quad \phi^{4} - 4\phi^{3} + 4\phi^{2} - 4\phi - 1 = 0;
\]
this equation may be written $2\phi(\phi - 2)(\phi^{2} + 1) = (\phi^{2} - 1)^{2}$, and we thence obtain
\[
\frac{(\phi^{2} - 1)^{4}}{16\phi^{2}(\phi^{2} + 1)^{2}} \;=\; \tfrac{1}{4}(\phi - 2)^{2};
\]
or the equation of the curve is $(x^{2} - 1)^{2} + (y^{2} - 1)^{2} = 1 + \tfrac{1}{4}(\phi - 2)^{2}$, where $\phi$ is determined by the equation just referred to. The curve is in this case symmetrical in regard to the two axes; and there are in fact four triangles, each in-and-circumscribed to the curve.

\medskip
\noindent\textit{Cambridge, June 16, 1865.}

\bigskip

% =====================================================================
% Paper 368 -- printed pp. 493-494, PDF pp. 519-520
% =====================================================================
\section*{Paper 368: On a Problem of Geometrical Permutation}

\medskip\noindent
[\textit{From the Philosophical Magazine, vol.\ \textsc{xxx}.\ (1865), pp.\ 370--372.}]

\medskip
It is required to find in how many modes the nine points of inflexion of a cubic curve can be denoted by the figures $1$, $2$, $3$, $4$, $5$, $6$, $7$, $8$, $9$, in such wise that the twelve lines, each containing three points of inflexion, shall be in every case denoted by the same triads of figures, say by the triads
\[
123,\; 147,\; 159,\; 168,
\]
\[
456,\; 258,\; 267,\; 249,
\]
\[
789,\; 369,\; 348,\; 357.
\]

We may imagine the inflexions so denoted in one particular way, which may be called the primitive denotation; then in any other mode of denotation, a figure, for example 1, is either affixed to the inflexion to which it originally belonged, and it is then said to be \emph{in loco}, or it is affixed to some other point of inflexion. This being so, the total number of modes is $=432$; viz.\ this number is made up as follows:
\[
\begin{array}{rl}
9 \text{ figures } in\;loco & 1\\
3 \;\; ,, \;\; ,, & 60\\
1 \text{ figure } & 243\\
0 \;\; ,, \;\; & 128\\\hline
& 432
\end{array}
\]

There is of course only one mode wherein the nine figures remain \emph{in loco}. It may be seen without much difficulty that there is not any mode in which $8$, $7$, $6$, $5$, or $4$ figures remain \emph{in loco}. There is no mode in which only $2$ figures remain \emph{in loco};

\newpage

\noindent for any two inflexions are in a line with a third inflexion; and if the figures which belong to the first two inflexions are \emph{in loco}, then the figure belonging to the third inflexion will be \emph{in loco}; that is, there will be 3 figures \emph{in loco}. The only remaining modes are therefore those which have 3 figures, 1 figure, or 0 figure \emph{in loco}.

First, if three figures are \emph{in loco}, these, as just seen, will be the figures which belong to three inflexions in a line. Suppose the figures are $1$, $2$, $3$; then the inflexion originally denoted, say by the figure $4$, may be denoted by any one of the remaining figures $5$, $6$, $7$, $8$, $9$; but when the figure is once fixed upon, then the remaining inflexions can be denoted only in one manner. Hence when the figures $1$, $2$, $3$ remain \emph{in loco} there are 5 modes; and consequently the number of modes wherein 3 figures remain \emph{in loco} is $5 \times 12$, $=60$.

Next, if only a single figure, suppose $1$, remains \emph{in loco}, the triads which belong to the figure $1$ are $123$, $147$, $159$, $168$; and there is 1 mode in which we simultaneously interchange all the pairs $(2,3)$, $(4,7)$, $(5,9)$, $(6,8)$. (Observe that the triads $123$, $147$, $159$, $168$ here denote the same lines respectively as in the primitive denotation, the figure $1$ remains \emph{in loco}, but the figures belonging to the other two inflexions on each of the four lines are interchanged.) There are, besides this, 2 modes in which the figures $(2,3)$, but not any other two figures, are interchanged; similarly 2 modes in which the figures $(4,7)$, 2 modes in which the figures $(5,9)$, 2 modes in which the figures $(6,8)$, but in each case no other two figures, are interchanged; this gives in all $1+2+2+2+2$, $=9$ modes. There are besides, the figure $1$ still remaining \emph{in loco}, 18 modes where there are no two figures $(2,3)$, $(4,7)$, $(5,9)$, or $(6,8)$ which are interchanged: viz.\ the figure $2$ may be made to denote any one of the inflexions originally denoted by $4$, $5$, $6$, $7$, $8$, or $9$. Suppose the inflexion originally denoted by $4$; $3$ will then denote the inflexion originally denoted by $7$: it will be found that of three of the remaining six inflexions, any one may be denoted by the figure $4$, and that the scheme of denotation can then in each case be completed in one way only. This gives $6 \times 3$, $=18$, as above, for the number of the modes in question; and we have then $9 + 18$, $=27$, for the number of the modes in which the figure $1$ remains \emph{in loco}; and $9 \times 27$, $=243$, for the number of modes in which some one figure remains \emph{in loco}.

Finally, if no figure remains \emph{in loco}, the figure $1$ will then denote some one of the inflexions originally denoted by $2$, $3$, $4$, $5$, $6$, $7$, $8$, $9$. Suppose it to denote that originally denoted by $2$; $2$ cannot then denote the inflexion originally denoted by $1$, for if it did, $3$ would remain \emph{in loco}: $2$ must therefore denote the inflexion originally denoted by $3$, or else some one of the inflexions originally denoted by $4$, $5$, $6$, $7$, $8$, $9$. It appears, on examination, that in the first case there are 4 ways of completing the scheme, and in each of the latter cases 2 ways; there are therefore in all $1 \times 4 + 6 \times 2$, $=16$ ways; that is, 16 modes in which (no figure remaining \emph{in loco}) the figure $1$ is used to denote the inflexion originally denoted by $2$; and therefore $8 \times 16$, $=128$ modes, for which no figure remains \emph{in loco}. This completes the investigation of the numbers $1$, $60$, $243$, and $128$, which together make up the total number $432$ of the modes of denotation of the nine inflexions.

\newpage

% =====================================================================
% Paper 369 -- printed pp. 495-497, PDF pp. 521-523
% =====================================================================
\section*{Paper 369: On a Property of Commutants}

\medskip\noindent
[\textit{From the Philosophical Magazine, vol.\ \textsc{xxx}.\ (1865), pp.\ 411--413.}]

\medskip
I call to mind the definition of a commutant, viz.\ if in the symbol
\[
\left[\begin{array}{c} 1\,1\,1\dots(\theta) \\ 2\,2\,2 \\ \vdots \\ p\,p\,p \end{array}\right]
\]
we permute independently in every possible manner the numbers $1, 2, \dots, p$ of each of the $\theta$ columns except the column marked $(\dagger)$, giving to each permutation its proper sign, $+$ or $-$, according as the number of inversions is even or odd, thus
\[
\pm A_{rst\dots(\theta)}
\]
which is to be read as meaning
\[
\pm A_{1r,1s,1t,\dots} \;A_{2r,2s,2t,\dots} \;\dots\; A_{pr,ps,pt,\dots};
\]
the sum of all the $(1\cdot 2\cdot 3 \dots p)^{\theta-1}$ terms so obtained is the commutant denoted by the above-mentioned symbol. In the particular case $\theta = 2$, the commutant is of course a determinant: in this case, and generally if $\theta$ be even, it is immaterial which of the columns is left unpermuted, so that the $(\dagger)$ instead of being placed over any column may be placed on the left hand of the $A$; but when $\theta$ is odd, the function has different values according as one or another column is left unpermuted, and the position of the $(\dagger)$ is therefore material. It may be added that if all the columns are permuted, then, if $\theta$ be even, the sum is $1 \cdot 2 \dots p$ into the commutant obtained by leaving any one column unpermuted; but if $\theta$ is odd, then the sum is $=0$.

\newpage

\noindent The property in question is a generalization of a property of determinants, viz.\ we have
\[
\begin{vmatrix}
2\lambda\lambda' & \lambda\mu' + \lambda'\mu & \lambda\nu' + \lambda'\nu & \dots\\
\lambda\mu' + \lambda'\mu & 2\mu\mu' & \mu\nu' + \mu'\nu & \dots\\
\lambda\nu' + \lambda'\nu & \mu\nu' + \mu'\nu & 2\nu\nu' & \dots\\
\vdots & \vdots & \vdots & \ddots
\end{vmatrix} = 0
\]
whenever the order of the determinant is greater than 2.

To enunciate the corresponding property of commutants, let
\[
\begin{pmatrix} \lambda_{11}, & \lambda_{12}, & \dots\\ \lambda_{21}, & \lambda_{22}, & \dots\\ \vdots & \vdots & \ddots \end{pmatrix}
\]
or, in a notation analogous to that of a commutant,
\[
\left[\begin{array}{c} \dagger\,1\,1 \\ \dagger\,2\,2 \\ \vdots\\ \dagger\,p\,p \end{array}\right]
\]
denote a function formed precisely in the manner of a determinant (or commutant of two columns), except that the several terms (instead of being taken with a sign $+$ or $-$ as above) are taken with the sign $+$: thus
\[
\left|\begin{array}{cc} \lambda_{11} & \lambda_{12}\\ \lambda_{21} & \lambda_{22} \end{array}\right|^{+} \quad\text{or}\quad \left[\begin{array}{c} \dagger\,1\,1\\ \dagger\,2\,2 \end{array}\right]
\]
each denote
\[
\lambda_{11}\lambda_{22} + \lambda_{12}\lambda_{21}.
\]
This being so, the theorem is that the commutant
\[
\left[\begin{array}{c} \dagger\,1\,1\,1\dots(\theta)\\ 2\,2\,2\\ \vdots\\ p\,p\,p \end{array}\right]
\]
where
\[
A_{rst\dots(\theta)} = \begin{pmatrix} \lambda_{1r}, & \lambda_{1s}, & \dots\,(\theta)\\ \lambda_{2r}, & \lambda_{2s}, & \dots\\ \vdots & \vdots & \ddots\\ \lambda_{pr}, & \lambda_{ps}, & \dots \end{pmatrix},
\]
whenever $p > \theta$, is $=0$.

To prove this, consider the general term of the commutant, viz.\ this is
\[
\pm A_{1's't'\dots(\theta)}\; A_{2's''t''\dots}\; \dots\; A_{p's^{(p)}t^{(p)}\dots};
\]

\newpage

\noindent the general term of $A_{rst\dots}$ is $\lambda_{ar}\lambda_{bs}\lambda_{ct}\dots$, where $a, b, c, \dots$ represent some permutation of the numbers $1, 2, 3, \dots, \theta$. Substituting the like values for each of the factors $A_{1's't'\dots}$, $A_{2's''t''\dots}$, etc., the general term of the commutant is
\[
= \pm \sum \lambda_{a'1}\lambda_{b's'}\lambda_{c't'}\dots\;\;\lambda_{a''2}\lambda_{b''s''}\lambda_{c''t''}\dots\;\dots\;\lambda_{a^{(p)}p}\lambda_{b^{(p)}s^{(p)}}\lambda_{c^{(p)}t^{(p)}}\dots.
\]
Taking the sum of this term with respect to the quantities $s', s'', \dots, s^{(p)}$, which denote any possible permutation of the numbers $1, 2, \dots, p$; again, with respect to the quantities $t', t'', \dots, t^{(p)}$, which denote any possible permutation of the numbers $1, 2, \dots, p$; and the like for each of the $(\theta - 1)$ series of quantities, the sum in question is
\[
\lambda_{a'1}\lambda_{a''2}\dots\lambda_{a^{(p)}p}\,\sum^{+}\!\lambda_{b'1}\lambda_{b''2}\dots\lambda_{b^{(p)}p}\,\sum^{+}\!\lambda_{c'1}\lambda_{c''2}\dots\lambda_{c^{(p)}p}\,\dots,
\]
which is
\[
\lambda_{a'1}\lambda_{a''2}\dots\lambda_{a^{(p)}p}\,\begin{pmatrix}\lambda_{b'1}\\ \lambda_{b''2}\\ \vdots\\ \lambda_{b^{(p)}p}\end{pmatrix}^{+}\!\begin{pmatrix}\lambda_{c'1}\\ \lambda_{c''2}\\ \vdots\\ \lambda_{c^{(p)}p}\end{pmatrix}^{+}\dots;
\]
but $p$ being greater than $\theta$, since the numbers $b', b'', \dots, b^{(p)}$ are all of them taken out of the series $1, 2, \dots, \theta$, some of these numbers must necessarily be equal to each other, and we have therefore
\[
\begin{pmatrix}\lambda_{b'1}\\ \lambda_{b''2}\\ \vdots\\ \lambda_{b^{(p)}p}\end{pmatrix}^{+} = 0,
\]
whence finally the commutant is $=0$.

In the case where $p = \theta = 2$, we have for a determinant of the order 2 the theorem
\[
\begin{vmatrix} 2\lambda\lambda' & \lambda\mu' + \lambda'\mu\\ \lambda\mu' + \lambda'\mu & 2\mu\mu' \end{vmatrix} = -\begin{vmatrix} \lambda & \mu\\ \lambda' & \mu' \end{vmatrix}^{2},
\]
and it is probable that there exists a corresponding theorem for the commutant
\[
\left[\begin{array}{c} \dagger\,1\,1\,1\dots(p)\\ 2\,2\,2\\ \vdots\\ p\,p\,p \end{array}\right],
\]
where
\[
A_{rst\dots(p)} = \begin{pmatrix} \lambda_{1r}, & \lambda_{1s}, & \dots\,(p)\\ \lambda_{2r}, & \lambda_{2s}, & \dots\\ \vdots & \vdots & \ddots\\ \lambda_{pr}, & \lambda_{ps}, & \dots \end{pmatrix},
\]
but I have not ascertained what this theorem is.

\medskip
\noindent\textit{Cambridge, October 26, 1865.}

\hfill 63

\bigskip

% =====================================================================
% Paper 370 -- printed pp. 498-499, PDF pp. 524-525
% =====================================================================
\section*{Paper 370: On the Signification of an Elementary Formula of Solid Geometry}

\medskip\noindent
[\textit{From the Philosophical Magazine, vol.\ \textsc{xxx}.\ (1865), pp.\ 413, 414.}]

\medskip
The expression for the perpendicular distance of a point $(x, y, z)$ from a line through the origin inclined at the angles $(\alpha, \beta, \gamma)$ to the three axes respectively, is
\begin{align*}
p^{2} &= x^{2} + y^{2} + z^{2} - (x\cos\alpha + y\cos\beta + z\cos\gamma)^{2}\\
&= (y\cos\gamma - z\cos\beta)^{2}\\
&\quad + (z\cos\alpha - x\cos\gamma)^{2}\\
&\quad + (x\cos\beta - y\cos\alpha)^{2};
\end{align*}
and the remark in reference to it is that, if at the given point $P$ we draw, perpendicular to the plane through $P$ and the given line, a distance $PK$ equal to the distance of $P$ from the given line, then the expressions
\[
y\cos\gamma - z\cos\beta,\quad z\cos\alpha - x\cos\gamma,\quad x\cos\beta - y\cos\alpha,
\]
which enter into the preceding formula, denote respectively the coordinates of the point $K$ referred to $P$ as origin.

If the given line instead of passing through the origin pass through the point $x_{0}, y_{0}, z_{0}$, then the corresponding expressions are of course
\[
(y - y_{0})\cos\gamma - (z - z_{0})\cos\beta,\quad (z - z_{0})\cos\alpha - (x - x_{0})\cos\gamma,\quad (x - x_{0})\cos\beta - (y - y_{0})\cos\alpha,
\]
and if we denote the ``six coordinates'' of the given line, viz.\
\[
\cos\alpha,\; \cos\beta,\; \cos\gamma,\; y_{0}\cos\gamma - z_{0}\cos\beta,\; z_{0}\cos\alpha - x_{0}\cos\gamma,\; x_{0}\cos\beta - y_{0}\cos\alpha,
\]
\[
a,\; b,\; c,\; f,\; g,\; h
\]
respectively (so that $af + bg + ch = 0$), then the three expressions become
\[
cy - bz - f,\quad az - cx - g,\quad bx - ay - h
\]
respectively.

It is moreover clear that if the point $P$ be moved to $P'$ by an infinitesimal rotation $\omega$ about the given line, then $P'$ lies on the line $PK$ at a distance $PP' = \omega\, PK$ from the point $P$, and the displacements of $P$ in the directions of the axes are consequently equal to
\[
\omega(cy - bz - f),\quad \omega(az - cx - g),\quad \omega(bx - ay - h)
\]
respectively, which is a fundamental formula in the theory of the infinitesimal rotations of a solid body.

\medskip
\noindent\textit{Cambridge, October 26, 1865.}

\hfill 63--2

\end{document}
